% !TEX encoding = UTF-8
% !TEX program = latexmk
\documentclass{CWNUExam}

\cwnuexamsetup{
pdf = {
title    = {计算机学院软件工程专业2024级JAVA编程语言期末考试试卷B卷},
author   = {Computer Science Department},
subject  = {For internal academic use only},
keywords = {Java; Programming; Final Exam},
},
}

\cwnuexamsetup{
title = {
year       = {\the\year},
semester   = {1},
course     = {JAVA Programming Language},
suffix     = {期末},
type       = {B},
total-part = {4},
college    = {计算机学院},
major      = {软件工程专业},
grade      = {2024级},
exam-type  = {闭卷},
exam-time  = {120},
exam-category = {考试},
test-type  = {normal},
},
}

\examsetup{page/size = a3paper}

\begin{document}
\maketitle

\section{Section A: Multiple Choice Questions (15 \times 3 = 45)}

\begin{question}
Which keyword is used to create a class in Java?
\paren
\begin{choices}
\item new
\item class
\item object
\item interface
\end{choices}
\end{question}

\begin{question}
Which data type is used to store decimal numbers?
\paren
\begin{choices}
\item int
\item char
\item float
\item boolean
\end{choices}
\end{question}

\begin{question}
Default value of a boolean variable is:
\paren
\begin{choices}
\item true
\item false
\item 0
\item null
\end{choices}
\end{question}

\begin{question}
Which symbol is used to access class members?
\paren
\begin{choices}
\item \texttt{:}
\item \texttt{.}
\item \texttt{\#} % 使用 \# 避免参数字符错误
\item \texttt{->}
\end{choices}
\end{question}

\begin{question}
Which loop is best used when the number of iterations is known?
\paren
\begin{choices}
\item for
\item while
\item do-while
\item foreach
\end{choices}
\end{question}

\begin{question}
Which keyword is used for inheritance in Java?
\paren
\begin{choices}
\item implements
\item extends
\item inherit
\item super
\end{choices}
\end{question}

\begin{question}
Method overriding occurs in:
\paren
\begin{choices}
\item same class
\item different classes
\item interface only
\item main method
\end{choices}
\end{question}

\begin{question}
Which access modifier allows access within the same package?
\paren
\begin{choices}
\item public
\item private
\item protected
\item default
\end{choices}
\end{question}

\begin{question}
Which keyword refers to the parent class object?
\paren
\begin{choices}
\item this
\item super
\item parent
\item base
\end{choices}
\end{question}

\begin{question}
Which of the following is an interface keyword?
\paren
\begin{choices}
\item class
\item interface
\item abstract
\item extends
\end{choices}
\end{question}

\begin{question}
Which exception occurs when dividing by zero?
\paren
\begin{choices}
\item IOException
\item ArithmeticException
\item NullPointerException
\item ArrayIndexOutOfBoundsException
\end{choices}
\end{question}

\begin{question}
Which block is used to handle exceptions?
\paren
\begin{choices}
\item try
\item catch
\item finally
\item all of these
\end{choices}
\end{question}

\begin{question}
Which feature of OOP hides data from outside access?
\paren
\begin{choices}
\item Inheritance
\item Polymorphism
\item Encapsulation
\item Abstraction
\end{choices}
\end{question}

\begin{question}
Which class is the parent of all Java classes?
\paren
\begin{choices}
\item System
\item Class
\item Object
\item Main
\end{choices}
\end{question}

\begin{question}
Which collection does NOT allow duplicate elements?
\paren
\begin{choices}
\item List
\item ArrayList
\item Set
\item Vector
\end{choices}
\end{question}



\section{Section B: Fill in the Blanks (5 \times 1 = 5)}

\begin{question}
The keyword used to define a class is \fillin.
\end{question}

\begin{question}
An array stores \fillin type of values.
\end{question}

\begin{question}
The keyword used to handle exceptions is \fillin.
\end{question}

\begin{question}
A method with the same name but different parameters is called \fillin.
\end{question}

\begin{question}
The \fillin block always executes whether an exception occurs or not.
\end{question}



% =========================
% Part C
% =========================
\section{Section C: Explain / Short Answer Questions (5 \times 2 = 10)}

\begin{problem}
What is a class in Java?
\end{problem}

\begin{problem}
What is an object?
\end{problem}

\begin{problem}
Explain inheritance in Java.
\end{problem}

\begin{problem}
What is polymorphism?
\end{problem}

\begin{problem}
What is exception handling?
\end{problem}


% =========================
% Part D
% =========================
\section{Section D: Programming / Design Questions (2 \times 20 = 40)}

\begin{problem}
\textbf{Design Question 1 (Student Management System)}

Design a simple Student Management system using inheritance.

\textbf{Guidelines:}
\begin{enumerate}
  \item Create a base class \texttt{Student} with: \texttt{name}, \texttt{rollNo}.
  \item Create a subclass \texttt{Result} that stores marks of three subjects.
  \item Add a method to calculate total marks.
  \item Display student details and result.
\end{enumerate}

\textbf{Note:}
\begin{itemize}
  \item Full Java code is not required.
  \item Class structure and key methods are sufficient.
\end{itemize}
\end{problem}

\begin{problem}
\textbf{Design Question 2 (Simple Bank System)}

Design a simple Bank system using interface.

\textbf{Guidelines:}
\begin{enumerate}
  \item Create an interface \texttt{Bank} with methods \texttt{deposit()} and \texttt{withdraw()}.
  \item Create a class \texttt{MyBank} that implements the interface.
  \item Use \texttt{try-catch} to handle a case where withdrawal amount is greater than balance.
\end{enumerate}

\textbf{Note:}
\begin{itemize}
  \item Pseudo-code or partial code is acceptable.
  \item Focus on program logic.
\end{itemize}
\end{problem}

        

% =========================
% Answer Sheet (One Page)
% =========================
\newpage

\noindent\textbf{Student Name:} \underline{\hspace{6cm}}

\textbf{Student ID:} \underline{\hspace{6cm}}

\vspace{0.6cm}
\noindent\textbf{Section A: Multiple Choice (1--15)}\\[0.2cm]
\noindent
\begin{tabular}{|c|c|c|c|c|}
\hline
1 & 2 & 3 & 4 & 5  \\
\hline
\underline{\hspace{1.6cm}} & \underline{\hspace{1.6cm}} & \underline{\hspace{1.6cm}} & \underline{\hspace{1.6cm}} & \underline{\hspace{1.6cm}}  \\
\hline
\end{tabular}

\begin{tabular}{|c|c|c|c|c|}
\hline
6 & 7 & 8 & 9 & 10  \\
\hline
\underline{\hspace{1.6cm}} & \underline{\hspace{1.6cm}} & \underline{\hspace{1.6cm}} & \underline{\hspace{1.6cm}} & \underline{\hspace{1.6cm}}  \\
\hline
\end{tabular}

\begin{tabular}{|c|c|c|c|c|}
\hline
11 & 12 & 13 & 14 & 15 \\
\hline
\underline{\hspace{1.6cm}} & \underline{\hspace{1.6cm}} & \underline{\hspace{1.6cm}} & \underline{\hspace{1.6cm}} & \underline{\hspace{1.6cm}}  \\
\hline
\end{tabular}

\vspace{0.8cm}
\noindent\textbf{Section B: Fill in the Blanks (1--5)}\\[0.2cm]
\noindent
1.\ \underline{\hspace{17cm}}\\[0.4cm]
2.\ \underline{\hspace{17cm}}\\[0.4cm]
3.\ \underline{\hspace{17cm}}\\[0.4cm]
4.\ \underline{\hspace{17cm}}\\[0.4cm]
5.\ \underline{\hspace{17cm}}\\[0.4cm]

\vspace{0.6cm}
\noindent\textbf{Section C: Short Answers (1--5)}\\[0.2cm]
\noindent
1.\ \underline{\hspace{17cm}}\\[0.35cm]
\underline{\hspace{17cm}}\\[0.5cm]
\underline{\hspace{17cm}}\\[0.5cm]
\underline{\hspace{17cm}}\\[0.5cm]
2.\ \underline{\hspace{17cm}}\\[0.35cm]
\underline{\hspace{17cm}}\\[0.5cm]
\underline{\hspace{17cm}}\\[0.5cm]
\underline{\hspace{17cm}}\\[0.5cm]
3.\ \underline{\hspace{17cm}}\\[0.35cm]
\underline{\hspace{17cm}}\\[0.5cm]
\underline{\hspace{17cm}}\\[0.5cm]
\underline{\hspace{17cm}}\\[0.5cm]
4.\ \underline{\hspace{17cm}}\\[0.35cm]
\underline{\hspace{17cm}}\\[0.5cm]
\underline{\hspace{17cm}}\\[0.5cm]
\underline{\hspace{17cm}}\\[0.5cm]
5.\ \underline{\hspace{17cm}}\\[0.35cm]
\underline{\hspace{17cm}}\\[0.5cm]
\underline{\hspace{17cm}}\\[0.5cm]
\underline{\hspace{17cm}}\\[0.5cm]
\vspace{0.6cm}
\noindent\textbf{Section D: Design Questions (1--2)}\\[0.2cm]
\noindent
\textbf{D1:}\\[0.2cm]
    \underline{\hspace{17cm}}\\[0.35cm]
\underline{\hspace{17cm}}\\[0.35cm]
\underline{\hspace{17cm}}\\[0.35cm]
\underline{\hspace{17cm}}\\[0.35cm]
\underline{\hspace{17cm}}\\[0.35cm]
\underline{\hspace{17cm}}\\[0.35cm]
\underline{\hspace{17cm}}\\[0.35cm]
\underline{\hspace{17cm}}\\[0.35cm]
\underline{\hspace{17cm}}\\[0.35cm]

\newpage
\textbf{D2:}\\[0.2cm]
\underline{\hspace{17cm}}\\[0.35cm]
\underline{\hspace{17cm}}\\[0.35cm]
\underline{\hspace{17cm}}\\[0.35cm]
\underline{\hspace{17cm}}\\[0.35cm]
\underline{\hspace{17cm}}\\[0.35cm]
\underline{\hspace{17cm}}\\[0.35cm]
\underline{\hspace{17cm}}\\[0.35cm]
\underline{\hspace{17cm}}\\[0.35cm]
\underline{\hspace{17cm}}\\[0.35cm]

\end{document}
